% Generated by scripts/assemble.py; corrections belong in source/corrections.json.
\examhead{55}{南开大学 2024 年数学分析试卷}{56}
% p057-b01; source PDF page 57
\par\Needspace{6\baselineskip}
\blockanchor{p057-b01}
\question{1}(20 分) 求极限\par
\nopagebreak[4]
% p057-b02; source PDF page 57
\blockanchor{p057-b02}
\displayformula{\lim _ {x \rightarrow 0} \left[ \frac {e ^ {x}}{\ln (1 + x)} - \frac {1}{x} \right].}
% p057-b03; source PDF page 57
\par\Needspace{6\baselineskip}
\blockanchor{p057-b03}
\question{2}(20 分) 设 $z$ 是由方程\par
\nopagebreak[4]
% p057-b04; source PDF page 57
\blockanchor{p057-b04}
\displayformula{2 x ^ {3} z - 3 x ^ {2} y ^ {2} - 2 x y z ^ {2} - 4 y ^ {3} z + 4 = 0}
% p057-b05; source PDF page 57
\blockanchor{p057-b05}
确定的关于 $x, y$ 的隐函数, 在 $x = 2, y = 1, z = 2$ 处计算全微分 $\mathrm{d}z$.\par
% p057-b06; source PDF page 57
\blockanchor{p057-b06}
\question{3}(20 分) 判断广义积分 $\int_0^{+\infty} \frac{x \ln x}{x^3 + 1} \, \mathrm{d}x$ 的敛散性.\par
% p057-b07; source PDF page 57
\par\Needspace{6\baselineskip}
\blockanchor{p057-b07}
\question{4}(25 分) 计算曲线积分\par
\nopagebreak[4]
% p057-b08; source PDF page 57
\blockanchor{p057-b08}
\displayformula{\int_ {L} \left(x ^ {3} y + e ^ {z}\right) \mathrm{d} x + \left(x y ^ {3} + 5 x ^ {3} y ^ {2} + x e ^ {y}\right) \mathrm{d} y.}
% p057-b09; source PDF page 57
\blockanchor{p057-b09}
其中 L 为椭圆  $4x^{2} + 9y^{2} = 36$ ，取逆时针方向.\par
% p057-b10; source PDF page 57
\blockanchor{p057-b10}
\question{5}(25 分) 求幂级数 $\sum_{n=1}^{\infty} n(n+1)\left(\frac{x-1}{x}\right)^n$ 的收敛域及和函数.\par
% p057-b11; source PDF page 57
\par\Needspace{6\baselineskip}
\blockanchor{p057-b11}
\question{6}(15 分) 设  $f(x)$  在 [0,1] 上单调递增,  $f(0)=0$ ,  $f(1)=1$ , n 为正整数,  $\xi_{k}\in\left[\frac{k-1}{n},\frac{k}{n}\right]$ ,  $k=1,2,\cdots,n$ , 证明:\par
\nopagebreak[4]
% p057-b12; source PDF page 57
\blockanchor{p057-b12}
\displayformula{\left| \int_ {0} ^ {1} f (x) \mathrm{d} x - \frac {1}{n} \sum_ {k = 1} ^ {n} f (\xi_ {k}) \right| \leq \frac {1}{n}.}
% p057-b13; source PDF page 57
\blockanchor{p057-b13}
\question{7}(15 分) 设 $f(x)$ 在 $[a, b]$ 上连续, $f(a) < f(b)$, 证明: 存在 $[c, d] \subseteq [a, b]$, 使得 $f(x)$ 在 $[c, d]$ 上的最小值为 $f(a)$, 最大值为 $f(b)$.\par
% p057-b14; source PDF page 57
\par\Needspace{6\baselineskip}
\blockanchor{p057-b14}
\question{8}(10 分) 设 $\sum_{n=1}^{\infty} a_n$ 是一个收敛的正项级数, 令 $f(x) = \sum_{n=1}^{\infty} a_n |\sin nx|$, 已知 $f(x)$ 在 $(- \infty, +\infty)$ 上满足利普希茨条件, 即存在常数 $L > 0$, 使得对任意的 $x, y \in (-\infty, +\infty)$, 均有\par
\nopagebreak[4]
% p057-b15; source PDF page 57
\blockanchor{p057-b15}
\displayformula{| f (x) - f (y) | \leq L | x - y |.}
% p057-b16; source PDF page 57
\blockanchor{p057-b16}
证明：级数 $\sum_{n=1}^{\infty} na_n$ 收敛.\par
